\chapter*{前言}

有一句著名的话叫做“程序=算法+数据结构”。我们使用算法来描述解决问题的步骤，而数据结构提供了\textbf{实现步骤}的方法。在学习数据结构时，首先应当明确的事情即是它支持的\textbf{操作}：以什么样的数据为输入，允许什么样的修改，能够回答什么样的查询。反过来，面对一个问题时，我们也要能够根据需要的操作选择合适的数据结构。

一般来说，数据结构的功能越强大，其代价也就越大。这种代价可以是多方面的：时空复杂度、实现难度、常数等等。也许 OIer 首次接触到通用、强大的方法就是在学习数据结构之时。在惊叹于强力的通用方法的同时，不应忽视那些功能较弱但代价也更低的方法，更不能放弃对问题本身的分析。挖掘问题的性质以使用更精练的方法，忽略问题的性质而使用强力的通用方法，两者都是处理问题时的重要思想。